\subsection{Design}
\label{uvm:sec:design}

% High-level overview of {\uvmname}'s design.
% TODO: brief paragraph motivating the four pieces below
%   and how they fit together.

\subsubsection{User-Space Caching}
\label{uvm:sec:design:caching}

The first piece of {\uvmname} is a UVM-aware,
    user-space caching allocator that lifts the
    \texttt{cudaMallocManaged}/\texttt{cudaFree} traffic off the critical path.
Without it, the most natural way to plug UVM into PyTorch---%
    a \texttt{CUDAPluggableAllocator} that forwards every allocation
    to \texttt{cudaMallocManaged} and every deallocation to \texttt{cudaFree}---%
    runs $54\times$ slower than vanilla PyTorch
    and even $1.6\times$ slower than Hugging~Face's
    application-level KV-cache offloading
    on OPT-13B (\S\ref{uvm:sec:bg:gap}).
The reason is structural:
    PyTorch's native \texttt{CUDACachingAllocator}
    retains freed blocks in user space and amortizes
    \texttt{cudaMalloc} cost across many tensor lifetimes,
    whereas the pluggable wrapper bypasses that cache and exposes
    every PyTorch \texttt{malloc}/\texttt{free} to the CUDA driver.
{\uvmname} closes this gap by porting the caching discipline
    of the native allocator to managed memory,
    while adding the extra bookkeeping that UVM requires
    to remain safe under oversubscription.

\paragraphb{CUDACachingAllocator-UVM}
{\uvmname} introduces \texttt{CUDACachingAllocator-UVM},
    a drop-in caching allocator that sits between
    PyTorch tensors and the CUDA driver,
    and that allocates its backing storage with \texttt{cudaMallocManaged}
    instead of \texttt{cudaMalloc}.
Its block-management policy mirrors the native allocator:
    blocks obtained from CUDA are split, coalesced,
    and reused across tensor lifetimes,
    so a PyTorch \texttt{free} merely returns a block to the cache
    rather than calling \texttt{cudaFree}.
As a result, only a small fraction of PyTorch-level allocations
    translates into a \texttt{cudaMallocManaged} call,
    and freed blocks remain warm for fast reuse.
Because the underlying memory is managed by UVM,
    these cached blocks are still eligible for transparent CPU spilling
    when the GPU is oversubscribed,
    so {\uvmname} retains both the throughput of caching
    and the oversubscription guarantees of UVM.

\paragraphb{UVM-Specific Proactive Garbage Collection}
Naively reusing the native caching policy on top of UVM is, however, unsafe.
The native allocator relies on \texttt{cudaMalloc} returning OOM
    as a back-pressure signal:
    when CUDA refuses a new allocation,
    the allocator invokes its garbage collector,
    releases cached blocks via \texttt{cudaFree},
    and retries.
With \texttt{cudaMallocManaged}, no such signal ever arrives---%
    the call simply succeeds by spilling pages to host memory.
Consequently, a UVM-backed caching allocator
    keeps growing its reservation indefinitely:
    GPU-resident pages overflow into CPU memory,
    cached but unused blocks are never returned,
    and the host kernel eventually OOM-kills the process
    long before UVM's spilling mechanism can stabilize the working set.

To eliminate this failure mode,
    {\uvmname} replaces the missing OOM signal
    with an explicit, GPU-side capacity check
    inside the allocator's slow path.
Before satisfying an allocation that would require
    a fresh \texttt{cudaMallocManaged} of size $s$,
    the allocator tests whether
    $s + \texttt{gpu\_reserved}$ would exceed the physical GPU memory,
    where $\texttt{gpu\_reserved}$ is the total bytes
    the allocator has already obtained from CUDA.
If so, it triggers proactive garbage collection:
    it walks its free-block list,
    releases reclaimable blocks back to CUDA via \texttt{cudaFree},
    and retries the allocation.
This keeps the allocator's reservation bounded
    by the physical GPU capacity,
    so UVM is only ever asked to oversubscribe
    on \emph{live} working-set growth,
    not on stale cached blocks---%
    avoiding the runaway CPU-side allocation that breaks naive UVM caching.

\paragraphb{Outcome}
Together, user-space caching and UVM-specific proactive GC
    turn UVM from an unusable substrate into a competitive one.
On the non-oversubscribed regime,
    {\uvmname}'s allocator is within a small overhead of vanilla PyTorch,
    delivering $51\times$ speedup over the pluggable-allocator UVM baseline
    and $16\times$ speedup over Hugging~Face's KV-cache offloading,
    while transparently enabling oversubscription
    without any application-level changes.
What remains, and what motivates the rest of this section,
    is that caching on top of UVM also \emph{couples} the allocator's cache
    to UVM's eviction decisions,
    creating fragmentation patterns that pure GC cannot resolve
    (\S\ref{uvm:sec:design:discard}).

\subsubsection{Guided Eviction with UVMDiscard}
\label{uvm:sec:design:discard}

User-space caching plus proactive GC stabilizes {\uvmname}'s reservation
    only when freed blocks can be returned to CUDA whole.
In real PyTorch workloads, however,
    cached blocks are routinely \emph{split} to satisfy smaller allocations,
    and the resulting fragments cannot be released independently---%
    a single \texttt{cudaMallocManaged} region can only be returned to CUDA
    via a single \texttt{cudaFree} of the entire region.
This couples caching with fragmentation in a way
    that pure GC cannot resolve,
    and pushes UVM into precisely the eviction traffic that
    PyTorch-level information would otherwise let us avoid.

Consider a sequence that allocates two $40$\,GiB blocks,
    frees them, then allocates two $25$\,GiB tensors
    that each carve a $25$\,GiB-used / $15$\,GiB-free split out of a cached block.
A subsequent $30$\,GiB allocation cannot be served from the cache:
    the two $15$\,GiB free fragments are non-contiguous
    and live inside still-live $40$\,GiB \texttt{cudaMallocManaged} regions,
    so GC cannot return them to CUDA either.
The allocator must therefore obtain a fresh $30$\,GiB region from CUDA,
    which on a now-reserved-full GPU forces UVM
    to evict roughly $30$\,GiB of pages to host memory.
Worse, the driver picks evictees from
    \emph{all} resident pages---including the $25$\,GiB \emph{live} tensors
    that the application is actively using---%
    incurring both non-trivial eviction traffic
    and the risk of evicting the wrong data.

\paragraphb{Discarding Partially-Freed Blocks}
{\uvmname} addresses this fragmentation--eviction coupling
    by extending its slow path with a second tier of reclamation
    that operates at sub-block granularity.
Before issuing a new \texttt{cudaMallocManaged}
    when $\texttt{new\_size} + \texttt{gpu\_reserved} > \texttt{gpu\_total}$,
    the allocator (i)~\texttt{cudaFree}s independent (fully unused) cached blocks,
    as in proactive GC,
    and then (ii)~for any cached block that is split into used and free fragments,
    invokes \texttt{cudaMemDiscardBatchAsync} on the free fragments
    to mark them discardable to the UVM driver.
Discard requires no application change:
    {\uvmname} already tracks the exact byte-range layout of every cached block
    as part of its allocator metadata,
    so identifying the partially-freed sub-regions is a constant-time lookup.
Crucially, the live fragments of the same block are left untouched,
    so existing tensors retain their virtual-address mappings
    and continue to fault and migrate as usual.

\paragraphb{Higher-Priority, Copy-Free Eviction}
Once a region is marked discarded,
    the UVM driver treats it differently from a normal resident page in two ways.
First, discarded pages are eligible to be evicted
    \emph{without} a device-to-host copy:
    the driver simply unmaps them and reclaims the physical frames,
    skipping the PCIe (or NVLink-C2C) transfer that dominates
    eviction cost on PCIe platforms.
Second, discarded pages are evicted at a \emph{higher priority}
    than normal resident pages,
    so when the driver needs to free GPU memory for a new allocation
    it draws from the discarded pool first.
For the fragmentation example above,
    discarding the two $15$\,GiB free fragments lets the driver
    satisfy the $30$\,GiB allocation by reclaiming exactly those $30$\,GiB---%
    with no copy and without touching live tensor data.

\paragraphb{Guiding UVM Away from Useless Data}
The combined effect is that {\uvmname} converts
    PyTorch-level "I no longer need these bytes" information
    into a UVM-level eviction hint
    that the driver can act on at page-fault time,
    without copying or evicting useful data.
This shifts the dominant source of eviction traffic
    from data the application still cares about
    to data the application has already discarded,
    cutting both eviction volume and end-to-end latency.
Empirically, discard makes UVM
    $1.6\times$ faster than the {\uvmname} configuration without discard
    and $27\times$ faster than Hugging~Face's KV-cache offloading
    on OPT-13B,
    even though the underlying allocator,
    proactive GC,
    and prepopulation policy are otherwise identical.

\subsubsection{Thrashing Mitigation Exploration}
\label{uvm:sec:design:thrash}
By default, {\uvmname} treats every CPU-resident page as a candidate for migration:
    when the GPU touches a non-resident page,
    the UVM driver faults it in
    so that subsequent accesses hit local HBM.
This is the right policy at low oversubscription
    --- it maximizes the fraction of the working set that lives on the GPU ---
    but at high oversubscription it degenerates into thrashing,
    because every fault-in implicitly forces an eviction
    of some other page that may itself be touched again moments later.
UVMDiscard (\S\ref{uvm:sec:design:discard}) reduces the cost of those evictions
    but does not prevent the eviction--re-fault cycle itself,
    and we observe that even with discard enabled
    LLM-inference workloads on a discrete-PCIe GPU
    spend most of their wall time on this cycle once oversubscription is significant.

We explore two complementary mechanisms that attack the cycle directly,
    by changing \emph{when} a page should be migrated rather than
    how cheaply it can be evicted.

\paragraphb{Direct CPU-Memory Access}
Instead of migrating pages on every fault,
    {\uvmname} can advise the driver to leave them on host
    (\texttt{cudaMemAdvise(\ldots, cudaMemAdviseSetPreferredLocation, hostId)})
    and let the GPU access them in place over the CPU--GPU interconnect.
No eviction or migration is incurred,
    so the eviction--re-fault cycle is short-circuited entirely
    --- the GPU pays a higher per-access latency
    in exchange for eliminating page-fault traffic.
This is most useful when the alternative would have been a near-immediate eviction anyway,
    i.e., when oversubscription is high enough that
    every fault-in steals a page that is itself still hot,
    and when the CPU--GPU link has enough bandwidth
    to absorb in-place accesses without becoming the new bottleneck.
We evaluate this trade-off on H100 PCIe in \S\ref{uvm:sec:eval:design:direct}.

\paragraphb{Access-Counter--Based Migration}
Modern NVIDIA GPUs maintain per-page access counters for UVM-managed memory,
    and the driver can be configured to defer migration
    until a host-resident page has been touched
    above a threshold number of times from the device.
This filters out one-shot accesses
    --- which would otherwise pay the cost of a full migration only to be evicted shortly after ---
    while still promoting genuinely hot pages to the GPU.
At present, access counters are only exposed on Grace-based systems
    with $64$\,KiB UVM pages,
    so we evaluate this mechanism on GH200 in \S\ref{uvm:sec:eval:design:counters};
    extending it to discrete-PCIe H100 (with $4$\,KiB pages)
    is contingent on driver support that is not yet available.

\subsubsection{Extension to System-Allocated Memory (SAM)}
\label{uvm:sec:design:sam}
\texttt{cudaMallocManaged} is not the only way to back UVM.
\emph{System-Allocated Memory} (SAM) achieves the same UVM properties
    --- transparent data movement and unified CPU/GPU addressing ---
    without going through a dedicated CUDA allocator,
    by letting the GPU access ordinary host pages obtained via
    \texttt{malloc}/\texttt{free}.
The OS exposes these pages to the GPU
    either through ATS (Address Translation Services) on GH200
    or through Linux's HMM path on x86 / H100,
    and we port {\uvmname} to both.

Two implementation differences matter when migrating from
    a \texttt{cudaMallocManaged}-backed UVM allocator
    to a SAM-backed one.

\paragraphb{Allocation alignment}
\texttt{cudaMallocManaged} returns pointers aligned to a coarse boundary
    (e.g.\ $512$\,B on GH200),
    and downstream CUDA libraries such as cuBLAS rely on this alignment
    when issuing tensor-cores and TMA-style accesses.
\texttt{malloc} provides only the C-ABI minimum alignment,
    which is finer-grained but insufficient for these consumers.
{\uvmname}'s SAM allocator therefore uses \texttt{posix\_memalign}
    to request the same alignment that \texttt{cudaMallocManaged} would have produced,
    so existing kernels work unchanged.

\paragraphb{Free-time synchronization}
\texttt{cudaMallocManaged} and \texttt{cudaFree} are CUDA runtime calls
    that internally synchronize against in-flight kernels and copies
    on the allocation's stream(s),
    so an application can free a buffer while a kernel referencing it is still queued
    without observing a use-after-free.
Plain \texttt{malloc}/\texttt{free} provide no such guarantee;
    on the SAM path, freeing a host page that the GPU is still touching
    yields exactly that bug.
{\uvmname}'s SAM allocator therefore wraps \texttt{free} with an explicit
    CUDA-side synchronization of all in-flight work
    on the streams that may have referenced the allocation,
    matching the semantics of the \texttt{cudaFree} path.

We kept the CUDA-level hint including cudaMemAdvise, cudaMemPrefetchAsync the same for SAM-based UVM.
UVMDiscard is currently not yet supported for SAM-based UVM.

Evaluation results can be found in \S\ref{uvm:sec:eval:sam}.


% TODO: describe how the same design extends to SAM
% (ATS on GH200, HMM on H100), including any allocator
% changes required to interoperate with malloc-backed memory.
